\appendix{}

\chapter{Derivations}\label{app:derivations}
\chapter{Extended Results}\label{app:results}
\chapter{Reproducibility}\label{app:repro}

\section*{Compute environment}

\autoref{tab:compute} is the machine of record for every number in this thesis. The hardware block
is transcribed from \texttt{sinfo}, \texttt{scontrol show node} and \texttt{lscpu} run on the cluster
on 2026-09-09, cross-checked against the provenance header (\texttt{NODE}, \texttt{GPU INFO},
\texttt{GIT REV}) that every batch job in this project writes to its own log. The software block is
transcribed from the environment-verification job that certified the cluster installation, and from
the pinned dependency list.

\begin{table}[htpb]
  \caption[Compute environment]{The machine and software stack of record. All runs used a single
  node; every job was allocated one GPU.}
  \label{tab:compute}
  \centering
  \begin{tabular}{ll}
    \toprule
    Component & Configuration \\
    \midrule
    \multicolumn{2}{l}{\emph{hardware}}\\
    Node                    & single shared GPU node, institute cluster \\
    CPU                     & $2\times$ AMD EPYC 7282, 16 cores each \\
    Cores / threads         & 32 physical / 64 logical, 2 NUMA domains, 1.5--2.8\,GHz \\
    Cache                   & 512\,KiB L2 per core; 128\,MiB L3 total ($8\times16$\,MiB) \\
    Main memory             & 252\,GiB; no node-local scratch \\
    GPUs on node            & $8\times$ NVIDIA RTX A5000, 24\,GB each \\
    GPUs per job            & \textbf{1}, held exclusively (no oversubscription, no preemption) \\
    NVIDIA driver           & 530.41.03, unchanged over the project \\
    \midrule
    \multicolumn{2}{l}{\emph{scheduling}}\\
    Scheduler               & SLURM 21.08.5 \\
    Wall-clock limit        & 24\,h per job, partition-enforced \\
    Typical training job    & 1 GPU, 8 threads, 32\,GB, 3--24\,h \\
    Operating system        & Ubuntu, Linux 5.15.0 \\
    \midrule
    \multicolumn{2}{l}{\emph{software}}\\
    Python                  & 3.10 (Conda environment) \\
    PyTorch                 & 2.2.2, CUDA 12.1 build \\
    NumPy / SciPy           & 1.26.4 / 1.13.1 \\
    MuJoCo                  & 2.3.7, headless EGL rendering \\
    Projection solver       & SciPy SLSQP (default); CasADi 3.6.5 with IPOPT selectable \\
    ODE integration         & \texttt{torchdiffeq} 0.2.3 \\
    Diffusers / Transformers & 0.31.0 / 4.41.2 \\
    Gym / Gymnasium         & 0.26.2 / 0.29.1 \\
    \bottomrule
  \end{tabular}
\end{table}

Two qualifications belong with the table. The software block is a snapshot taken when the
environment was certified at the start of the project; no version change is recorded in any
subsequent job log, but the environment was not frozen, so the defensible claim is that these
versions were verified at the outset rather than that each one produced every reported number.
And the accelerator, though the GPU is held exclusively, sits in a node whose CPU and memory are
shared, which is why \autoref{sec:setup:protocol} declines to read the timing figures as a hardware
benchmark.

Reproducing a run requires the batch scripts as much as the code: the scheduling parameters, the
GPU pinning and the headless-rendering configuration all live in the job scripts, and each job
records the git revision it ran at in its own log header.
\srcnote{Slurm\_Codes/submit.sh; Slurm\_Codes/sbatch/}

\hole{Total compute spent is not yet quantified. Arithmetic over the batch-log timestamps gives an
order of $10^{3}$ GPU-hours across roughly five months, but that covers only the jobs whose logs
record both a start and an end; re-derive it with \texttt{sacct} before stating any figure.}

